\section{世界模型、触觉与灵巧手}

前面讨论动作迁移和数据，本章进入两个更底层的能力：世界模型与触觉。谭杰把世界模型描述为 Vision-Language-Vision：输入当前视觉和语言/动作，生成下一帧或未来视觉状态。触觉则在灵巧手场景里变得重要，因为精细抓取、滑动、材质和接触力无法只靠视觉稳定解决。

\begin{figure}[H]
\centering
\includegraphics[width=0.96\textwidth]{world-model-vlv.png}
\caption{世界模型 V-L-V：Vision + Language 输入，预测下一帧视觉状态。自制概念图，依据 01:03:48--01:08:29 对谈内容整理。}
\end{figure}

\begin{knowledgebox}{读图：世界模型是“预测动作后果”}
当前画面和语言目标进入 World Model，模型预测下一帧或未来视觉状态，Policy 再据此选择动作。它不是只生成漂亮视频，而是帮助机器人想象“如果我这么做，世界会怎样”。
\end{knowledgebox}

\subsection{世界模型为什么重要}

本节把世界模型从“视频生成”拉回机器人控制。机器人需要的是预测行动后果，而不是生成一段看起来合理的画面。

机器人需要在行动前预测后果。抓杯子会不会倒，拉链往哪个方向拉，手碰到布料会怎样变形，这些都需要对世界状态变化建模。世界模型如果足够强，可以减少真实试错成本，提高规划和恢复能力。

\subsection{触觉为什么重新重要}

接下来讨论传感器层面的变化。过去触觉常被认为不重要，部分原因是硬件能力没有到需要精细触觉的阶段；一旦进入灵巧手，触觉就从锦上添花变成基础输入。

谭杰提到，过去觉得触觉不重要，是受限于硬件；如果有灵巧手，触觉就非常重要。普通夹爪可以靠视觉和粗动作完成不少任务；灵巧手则需要知道接触力、滑动、材质、抓取稳定性和指尖状态。没有触觉，很多精细 manipulation 很难可靠完成。

\begin{figure}[H]
\centering
\includegraphics[width=0.94\textwidth]{tactile-dexterous-hand.png}
\caption{触觉与灵巧手：硬件越灵巧，触觉越重要。自制概念图，依据 01:08:29--01:17:35 对谈内容整理。}
\end{figure}

\begin{knowledgebox}{读图：触觉价值随本体复杂度上升}
普通夹爪主要处理位置和闭合；灵巧手要处理多指协同、接触力、滑动、材质和抓取稳定性。硬件越接近人的手，触觉越不再是可有可无的传感器。
\end{knowledgebox}

\subsection{本章小结}

世界模型解决“预测后果”，触觉解决“精细接触”。二者都指向同一件事：机器人要从粗糙动作走向可靠操作，必须理解物理世界的变化和接触。

